% Part:sets-functions-relations
% Chapter: size-of-sets
% Section: reduction-alt

\documentclass[../../../include/open-logic-section]{subfiles}

\begin{document}

\olfileid{sfr}{siz}{red-alt}

\olsection{ન્યૂનીકરણ}

\begin{editorial}
  આ વિભાગ \olref[nen-alt]{sec}નાં પરિણામોને અનુરૂપ રીતે
  ન્યૂનીકરણ દ્વારા અગણનીયતા સાબિત કરે છે. \olref[nen]{sec}નાં
  પરિણામોને અનુરૂપ
  સહેજ વધુ વિસ્તૃત વૈકલ્પિક સ્વરૂપ \olref[red]{sec}માં આપેલું છે.
\end{editorial}

વિકર્ણીકરણની દલીલ વડે આપણે $\Bin^\omega$ એ !!{nonenumerable}
છે તે સાબિત કર્યું. સમાન વિકર્ણીકરણની દલીલથી $\Pow{\Nat}$ એ
!!{nonenumerable} છે તે બતાવ્યું. પરંતુ $\Pow{\Nat}$ એ
!!{nonenumerable} છે તે સાબિત કરવાની
આ બીજી રીત છે:
\emph{જો $\Pow{\Nat}$ એ !!{enumerable} હોય, તો $\Bin^\omega$ પણ
!!{enumerable} છે} તે બતાવો. $\Bin^\omega$ એ !!{nonenumerable}
છે તે આપણે જાણીએ છીએ, તેથી $\Pow{\Nat}$ પણ અગણનીય છે તેમ મળશે.

આને એક સમસ્યાનું બીજી સમસ્યામાં \emph{ન્યૂનીકરણ} કહે છે. આ કિસ્સામાં,
આપણે $\Bin^\omega$ને પરિગણિત કરવાની સમસ્યાનું $\Pow{\Nat}$ને
પરિગણિત કરવાની સમસ્યામાં ન્યૂનીકરણ કરીએ છીએ. બીજી સમસ્યાનો
ઉકેલ---$\Pow{\Nat}$ની પરિગણના---પહેલી સમસ્યાનો
ઉકેલ---$\Bin^\omega$ની પરિગણના---આપશે.

ગણ~$B$ને પરિગણિત કરવાની સમસ્યાનું ગણ~$A$ને પરિગણિત કરવાની
સમસ્યામાં ન્યૂનીકરણ કરવા, આપણે $A$ની પરિગણનાને $B$ની પરિગણનામાં
ફેરવવાની રીત આપીએ છીએ. તે કરવાની સૌથી સરળ રીત !!a{surjection}
$f\colon A \to B$ વ્યાખ્યાયિત કરવાની છે. જો $x_1$, $x_2$, \dots{}
$A$ની પરિગણના કરે, તો $f(x_1)$, $f(x_2)$, \dots{} એ $B$ની
પરિગણના કરશે. આપણા કિસ્સામાં આપણે !!a{surjection}
$f\colon \Pow{\Nat} \to \Bin^\omega$ શોધીએ છીએ.

\begin{prob}
જો કોઈ !!{injective} વિધેય $g\colon B \to A$ હોય અને $B$ એ
!!{nonenumerable} હોય, તો $A$ પણ અગણનીય છે તે બતાવો. $A$ની
પરિગણનામાંથી $B$ની પરિગણના બનાવવા~$g$નો ઉપયોગ કેવી રીતે કરી
શકાય તે બતાવીને આ કરો.
\end{prob}

\begin{proof}[ન્યૂનીકરણ દ્વારા {\olref[nen-alt]{thm:nonenum-pownat}}ની સાબિતી]
ન્યૂનીકરણ માટે ધારો કે $\Pow{\Nat}$ એ !!{enumerable} છે અને તેથી
તેની કોઈ પરિગણના $N_{1}$, $N_{2}$, $N_{3}$, \dots છે.

વિધેય $f \colon \Pow{\Nat} \to \Bin^\omega$ને આ રીતે વ્યાખ્યાયિત
કરો: $f(N)$ એ એવી પ્રતીકશ્રેણી~$s$ છે કે $s(n) = 1$ ત્યારે અને
માત્ર ત્યારે જ જ્યારે $n \in N$, અને અન્યથા $s(n) = 0$.
\sourcecorrection{OLSIZ-010}{મૂળની \texttt{reduction-alt.tex},
પંક્તિઓ 53--56માં \(f(N)\)ની કિંમત માટે અનિર્ધારિત સૂચકવાળું
\(s_k\) લખાયું છે. પછીની વ્યાપ્તતાની દલીલ મનસ્વી \(s\) વાપરે છે;
અહીં આખી વ્યાખ્યામાં \(s\) સુસંગત રીતે વાપર્યું છે.}

આ સ્પષ્ટપણે વિધેય વ્યાખ્યાયિત કરે છે, કારણ કે જ્યારે પણ
$N \subseteq \Nat$ હોય, ત્યારે કોઈ પણ $n \in \Nat$ એ $N$નો
!!a{element} હોય અથવા ન હોય. દાખલા તરીકે, યુગ્મ પ્રાકૃતિક સંખ્યાઓનો
ગણ $2\Nat = \Setabs{2n}{n \in \Nat} = \{0,2, 4, 6,
\dots\}$નું નિરૂપણ પ્રતીકશ્રેણી $1010101\dots$ પર થાય છે;
$\emptyset$નું $0000\dots$ પર; અને $\Nat$નું $1111\dots$ પર થાય છે.

તે !!{surjective} પણ છે: $0$ અને $1$ની દરેક પ્રતીકશ્રેણી પ્રાકૃતિક
સંખ્યાઓના કોઈ ગણને અનુરૂપ છે, એટલે કે પ્રતીકશ્રેણીમાં જે સ્થાનો પર
$1$ હોય તેને અનુરૂપ પ્રાકૃતિક સંખ્યાઓ તેના સભ્યો છે. વધુ ચોક્કસ રીતે,
જો $s \in \Bin^\omega$ હોય, તો $N \subseteq \Nat$ને આ રીતે
વ્યાખ્યાયિત કરો:
\[
N = \Setabs{n \in \Nat}{s(n) = 1}
\]
ત્યારે $f$ની વ્યાખ્યા તપાસીને $f(N) = s$ છે તે ચકાસી શકાય છે.

હવે આ યાદી વિચારો:
\[
f(N_1), f(N_2), f(N_3), \dots
\]
$f$ એ !!{surjective} હોવાથી $\Bin^\omega$નો દરેક સભ્ય કોઈ
દલીલ માટેની~$f$ની કિંમત તરીકે આવવો જોઈએ અને તેથી યાદીમાં આવવો જ
જોઈએ. આથી આ યાદી આખા~$\Bin^\omega$ની પરિગણના કરે છે.

એટલે જો $\Pow{\Nat}$ એ !!{enumerable} હોત, તો $\Bin^\omega$
!!{enumerable} હોત. પરંતુ $\Bin^\omega$ એ !!{nonenumerable} છે
(\olref[nen-alt]{thm:nonenum-bin-omega}). તેથી $\Pow{\Nat}$ એ
!!{nonenumerable} છે.
\end{proof}

%\begin{explain}
%ન્યૂનીકરણ કઈ દિશામાં જાય છે તેમાં ગૂંચવાવું સહેલું છે.
%દાખલા તરીકે, !!a{surjective} વિધેય $g \colon \Bin^\omega \to X$
%હોવાથી $X$ એ !!{nonenumerable} છે એવું \emph{સ્થાપિત થતું નથી}.
%($g(s) = s(1)$ વડે વ્યાખ્યાયિત વિધેય $g
%\colon \Bin^\omega \to \Bin$નો વિચાર કરો, જે $0$ અને $1$ના
%અનુક્રમને તેના પહેલા !!{element} પર મોકલે છે. તે વ્યાપ્ત છે,
%કારણ કે કેટલાક અનુક્રમો $0$થી અને કેટલાક $1$થી શરૂ થાય છે.
%પરંતુ $\Bin$ સાન્ત છે.) વળી, વિધેય $f$ વ્યાપ્ત હોવું જ જોઈએ;
%અન્યથા દલીલ ચાલતી નથી: $f(x_1)$, $f(x_2)$, \dots{}માં
%$Y$ના બધા !!{element}s આવે તેની ખાતરી રહેતી નથી. દાખલા તરીકે,
%$h\colon \Nat \to
%\Bin^\omega$ને આ રીતે વ્યાખ્યાયિત કરો:
%\[
%h(n) = \underbrace{000\dots0}_{\text{$n$ $0$'s}}
%\]
%આ એક વિધેય છે, પરંતુ $\Nat$ એ !!{enumerable} છે.
%\end{explain}

\begin{prob}\label{sfr:siz:red-alt:prob:nat-nat}
\sourcecorrection{OLSIZ-012}{મૂળની \texttt{reduction-alt.tex},
  પંક્તિ 107માં આ વૈકલ્પિક વિભાગના પ્રશ્નને પ્રાથમિક
  \texttt{red} વિભાગનું પહેલેથી વપરાયેલું સંપૂર્ણ label અપાયું છે.
  અહીં તેને વર્તમાન \texttt{red-alt} વિભાગનું અનન્ય label આપ્યું છે.}
  ન્યૂનીકરણની દલીલ વડે બતાવો કે બધા વિધેયો
  $f\colon \Nat \to \Nat$નો ગણ~$X$ !!{nonenumerable} છે.
  (સંકેત: $X$થી~$\Bin^\omega$માં વ્યાપ્ત વિધેય આપો.)
\end{prob}

\begin{prob}
ન્યૂનીકરણની દલીલ વડે બતાવો કે પ્રાકૃતિક સંખ્યાઓની જોડોના બધા
\emph{ગણોનો} ગણ, એટલે કે $\Pow{\Nat \times \Nat}$, !!{nonenumerable} છે.
\end{prob}

\begin{prob}
ન્યૂનીકરણની દલીલ વડે બતાવો કે પ્રાકૃતિક સંખ્યાઓના અનંત
અનુક્રમોનો ગણ~$\Nat^\omega$ !!{nonenumerable} છે.
\end{prob}

%\begin{prob}
%$P$ એ $\Nat$થી ગણ $\{0\}$માંના વિધેયોનો ગણ અને $Q$ એ ધન
%પૂર્ણાંકોના ગણથી ગણ $\{0\}$માંના \emph{આંશિક} વિધેયોનો ગણ હોય.
%બતાવો કે $P$ એ !!{enumerable} છે અને $Q$ નથી. (સંકેત:
%$\Bin^\omega$ને પરિગણિત કરવાની સમસ્યાનું~$Q$ને પરિગણિત કરવાની
%સમસ્યામાં ન્યૂનીકરણ કરો.)
%\end{prob}

\begin{prob}
$S$ એ $\Nat$થી ગણ $\{0,1\}$માંના બધા !!{surjection}sનો ગણ હોય;
એટલે કે $S$માં બધા !!{surjection}s~$f \colon \Nat \to \Bin$ છે.
બતાવો કે $S$ એ !!{nonenumerable} છે.
\end{prob}

\begin{prob}
બતાવો કે બધી વાસ્તવિક સંખ્યાઓનો ગણ~$\Real$ !!{nonenumerable} છે.
\end{prob}

\end{document}
